\scanpage{204}
\mspage{L 17}
entrevu seulement et encore moins vu et touché, \deleted{pour} peu à peu se
\deleted{dégager de son manteau \inserted{\deleted{d'ombre et}} de brumes, et se}
décante\deleted{r} de l'impondérable, se dégage\deleted{r} de son manteau d'ombre et de brumes
\deleted{et}\inserted{pour} prendre forme et chair\inserted{et poids}\inserted{\ldots}
\margin{\figure{\begin{tikzpicture}[scale=.3]
\draw[blue] (0,0)--(1,0)--(1,1)--(2.5,1);\draw[blue] (.9,0)--(.9,.9);
\end{tikzpicture}}}

\inserted{\figure{\begin{tikzpicture}[scale=.25]\draw[blue] (0,0)--(.8,0)--(.8,1)--(2,1);\draw[blue] (.65,0)--(.65,1);\end{tikzpicture}}}\deleted{Ce sont}\inserted{C'est}\deleted{ces stades}\inserted{cette partie-}là du travail,
de piètre apparence pour ne pas dire\inserted{(bien des fois)} foireux, qui
\deleted{\unsure{en représentent aussi les points}}\inserted{en \deleted{sont} est aussi la \deleted{plus}}
\deleted{névralgique}\inserted{partie la plus délicate et la plus essentielle}
--- celle où, véritablement quelque chose de n\,o\,u\,v\,e\,a\,u fait son apparition,
par l'effet d'une attention intense, d'une sollicitude, d'un respect pour cette chose fragile,
infiniment délicate, sur le point de naître. C'est la partie créatrice entre toutes ---
celle\deleted{s} de la conception et d\deleted{e la}\inserted{'une}\inserted{lente} gestation
\deleted{d'un fragile embryon} dans les chaudes ténèbres de la matrice nourricière,
\deleted{\unsure{évoluant en silence}} depuis l'invisible double gamète \deleted{initiale}
originelle devenant\inserted{informe} embryon et se transformant au fil des jours et des mois,
par un travail obscur et intense, invisible et sans apparence, en un nouvel être en chair et en os.

C'est là aussi la partie\inserted{éminemment} ``yin'', ``féminine\deleted{\unsure{,}}'' du travail de découverte\inserted{.}\deleted{---}
\deleted{alors que \deleted{la partie}}\deleted{l'}\inserted{L'}aspect complémentaire ``yang'', ou masculin,
s'apparenterait plutôt \deleted{aux} au travail à coups de marteaux ou de masse,
sur un burin bien affûté ou sur un coin de bon acier trempé.
L'un et l'autre aspect a sa raison d'être et sa fonction, en symbiose inséparable\inserted{l'un}
avec l'autre --- ou pour mieux dire, ce sont là l'\uline{épouse} et \uline{l'époux}
du couple des deux forces cosmiques originelles, dont l'étreinte
\inserted{\deleted{sans cesse}}\deleted{renouvelée fait}
sans cesse renouvelée fait resurgir sans cesse les obscurs labeurs créateurs de la conception,
de la gestation et de la naissance --- de la naissance de l'\,e\,n\,f\,a\,n\,t, de la chose nouvelle.

La deuxième chose sur laquelle je sentais le besoin de m'exprimer, dans ma fameuse ``introduction''
\deleted{philosophique et} personnelle et ``philosophique'' à un texte mathématique,
c'était au sujet de la nature du travail créateur justement.
Je me rendais compte, depuis des années, que cette nature était généralement ignorée,
occultée par des clichés à tout venant et par des\inserted{répressions et des} peurs ancestrales.
A quel point il en est bien ainsi, je l'ai découvert\inserted{\unsure{après seulement,|alors seulement,}}
progressivement, en fil des jours et des mois, tout au cours de la réflexion et de ``l'enquête''
poursuivie dans Récoltes et Semailles.
C'est dès les tout débuts de cette réflexion, au mois de juin 1983, que je suis saisi pour la première fois
(pour l'oublier\inserted{un peu} par la suite\deleted{\ldots) par}
\inserted{, sous le flux des découvertes imprévues\ldots) par}
la portée de ce fait d'anodine apparence, et pourtant \unsure{asses|assez} stupéfiant,
pour peu seulement qu'on s'y arrête tant soit peu: que cette partie ``créatrice entre toutes''
dont je viens de parler dans le travail de découverte,
\uline{ne transparaît pratiquement nulle part} dans les textes\inserted{ou discours} qui sont censés
\margin{203}
